\documentclass[11pt]{article}

%──────────────────────────────────────────────────────────────────────────────%
%                               PACKAGES                                      %
%──────────────────────────────────────────────────────────────────────────────%
\usepackage[utf8]{inputenc}
\usepackage[T1]{fontenc}
\usepackage{tgschola} % Different font
\usepackage{geometry}
  \geometry{a4paper, margin=1in}
\usepackage{amsmath,amssymb,amsthm}
\usepackage{mathtools}
\usepackage{xcolor}
\usepackage{hyperref}
\usepackage{enumitem}
\usepackage{graphicx}
\usepackage{tikz}
\usepackage{comment}
\usepackage{mathrsfs}   % script fonts
\usepackage{bm}         % bold math (for vectors)
\usepackage[all,cmtip]{xy}
\usepackage{titlesec} % For section styling

%──────────────────────────────────────────────────────────────────────────────%
%                               COLOURS                                        %
%──────────────────────────────────────────────────────────────────────────────%
\definecolor{MyPurple}{RGB}{128, 0, 128}
\definecolor{MyDarkBlue}{RGB}{0, 0, 139}
\definecolor{MyDarkGreen}{RGB}{0, 100, 0}
\definecolor{MyRed}{RGB}{204, 0, 0}

\hypersetup{
    colorlinks=true,
    linkcolor=black,
    citecolor=black,
    urlcolor=black
}

%──────────────────────────────────────────────────────────────────────────────%
%                          TITLE & SECTION STYLING                             %
%──────────────────────────────────────────────────────────────────────────────%
\titleformat{\section}
  {\color{MyDarkGreen}\normalfont\Large\bfseries}
  {\thesection}{1em}{}
\titleformat{\subsection}
  {\color{MyPurple}\normalfont\large\bfseries}
  {\thesubsection}{1em}{}
\renewcommand{\paragraph}[1]{\vspace{1em}\noindent{\color{MyDarkBlue}\bfseries #1}}

%──────────────────────────────────────────────────────────────────────────────%
%                          THEOREM-LIKE ENVIRONMENTS                         %
%──────────────────────────────────────────────────────────────────────────────%

\newtheoremstyle{mytheoremstyle}% name
  {6pt}% Space above
  {6pt}% Space below
  {\itshape}% Body font
  {}% Indent amount
  {\color{MyPurple}\bfseries}% Theorem head font
  {.}% Punctuation after theorem head
  {0.5em}% Space after theorem head
  {}% Theorem head spec (can be left empty)

\theoremstyle{mytheoremstyle}
\newtheorem{theorem}{Theorem}[section]
\newtheorem{proposition}[theorem]{\color{MyPurple}Proposition}

\theoremstyle{definition}
\newtheorem{definition}[theorem]{\color{MyDarkBlue}Definition}
\newtheorem{example}[theorem]{\color{MyDarkGreen}Example}

\theoremstyle{remark}
\newtheorem{remark}[theorem]{\color{MyDarkGreen}Remark}
\newtheorem*{note}{\color{MyDarkGreen}Note}


%──────────────────────────────────────────────────────────────────────────────%
%                           CUSTOM MACROS & SHORTCUTS                          %
%──────────────────────────────────────────────────────────────────────────────%
\DeclareMathOperator{\Tr}{Tr}
\newcommand{\bbC}{\mathbb{C}}
\newcommand{\ot}{\otimes}
\newcommand{\cU}{\mathcal U}
\newcommand{\C}{\mathbb{C}}
\newcommand{\diag}{\operatorname{diag}}
\newcommand{\bbR}{\mathbb{R}}
\newcommand{\cA}{\mathcal{A}}
\newcommand{\cB}{\mathcal{B}}
\newcommand{\bbN}{\mathbb{N}}
\newcommand{\abs}[1]{\lvert#1\rvert}
\newcommand{\norm}[1]{\lVert#1\rVert}
\newcommand{\set}[1]{\{\,#1\,\}}
\newcommand{\cW}{\mathcal{W}}
\newcommand{\cZ}{\mathcal{Z}}
\newcommand{\Cst}{$C^{*}$}
\newcommand{\dimvec}[1]{\widetilde{\bm #1}} % dimension column-vector
\newcommand{\Hom}{\operatorname{Hom}}
\newcommand{\End}{\operatorname{End}}
\newcommand{\tr}{\operatorname{tr}}
\newcommand{\id}{\mathrm{id}}

\newcommand{\missingfigure}[1]{{\color{MyRed} \bfseries [Missing figure: #1]}}

%──────────────────────────────────────────────────────────────────────────────%
%                                 FRONT MATTER                                %
%──────────────────────────────────────────────────────────────────────────────%
\title{\color{MyDarkBlue}\bfseries\huge Note on UOB Project}
\author{Jay}
\date{\today}

\begin{document}
\maketitle

\begin{abstract}
\noindent
\textit{Here, we include all the process on this topic. We will start by putting the Primary example, and step by step go on.}
\end{abstract}

\tableofcontents
\vspace{1em}


\newpage
%──────────────────────────────────────────────────────────────────────────────%
%                                SAMPLE USAGE                                 %
%──────────────────────────────────────────────────────────────────────────────%
\section{Introduction}

\subsection*{Finite–dimensional \Cst-algebras}

\begin{definition}[{\cite[§3.2]{JonesSunder}}]
A \emph{finite–dimensional \Cst-algebra} $A$ is $\ast$-isomorphic to a finite
direct sum of full matrix algebras
\[
  A \cong \bigoplus_{i=1}^{s} M_{n_i}(\mathbb{C}) , \qquad (n_i\ge 1).
\]
\end{definition}

\paragraph{Structure theorem.}
Every finite–dimensional \Cst-algebra is semisimple; its simple components are
the matrix blocks above and its \textbf{minimal central projections}
\[
  p_i=(0,\dots,0,\,I_{n_i},\,0,\dots,0)
\]
sum to $1_A$.

\paragraph{Dimension vector.}
Write
\(
  \dimvec n = (n_1,\dots,n_s)^{\!\mathsf T}\!,\)
a \emph{column} vector.  For each $i$, the irreducible left $A$-module
$V_i=\mathbb{C}^{n_i}$ satisfies
$\dim_{\mathbb{C}}V_i=n_i$.

\begin{definition}
Let $B=\bigoplus_{j=1}^{r}M_{m_j}(\mathbb{C})$ be another finite–dimensional
\Cst-algebra with dimension vector
\[
  \dimvec m=(m_1,\dots,m_r)^{\!\mathsf T}.
\]
A \emph{unital inclusion} $\,B\subseteq A\,$ is specified by the
\emph{inclusion matrix}
\[
H=[h_{ij}]\in M_{s\times r}(\mathbb{N}_0)
\]
whose entry $h_{ij}$ is the multiplicity of the simple $B$-module
$W_j=\mathbb{C}^{m_j}$ inside $V_i\!\!\restriction_B$.  Equivalently
\begin{equation}\label{eq:Hmn}
  H\,\dimvec m =\dimvec n.
\end{equation}
\end{definition}

%--------------------------------------------------------------------
\subsection*{Conditional expectations and traces}

%--------------------------------------------------------------------
\begin{definition}[Conditional expectation]\label{def:cond-exp}
Let $A$ be a unital $C^{*}$-algebra and $B\subseteq A$ a unital (i.e. $1_{B}=1_{A}$)
$C^{*}$-subalgebra.  A \emph{conditional expectation} (of $A$ \emph{onto} $B$) is a linear map
\[
   E:A\longrightarrow B
\]
satisfying the following four properties:
\begin{enumerate}
  \item\textbf{Completely positive:} for every $n\in\mathbb{N}$ the
        amplified map
        \[
        E^{(n)}:M_{n}(A)\to M_{n}(B),\; [x_{ij}]\mapsto[E(x_{ij})]
        \]
sends positive matrices to positive matrices.

  \item\textbf{Unital (unit–preserving):} $E(1_{A})=1_{B}$.

  \item\textbf{$B$–bimodular:} 
        \[
        E(b_{1}x\,b_{2}) = b_{1}\,E(x)\,b_{2}
        \quad
        \text{for all }b_{1},b_{2}\in B,\;x\in A.
        \]

  \item\textbf{Idempotent (projection):} $E^{2}=E$
        (equivalently $E|_{B}=\mathrm{id}_{B}$). 
\end{enumerate}
\end{definition}


\paragraph{Canonical (\textbf{Markov}) trace.}
Write $\tr_{n_i}$ for the normalised matrix trace on $M_{n_i}(\mathbb{C})$.
The \emph{standard trace} on $A$ is
\[
  \tau_A\!\bigl((x_1,\dots,x_s)\bigr)
  =\sum_{i=1}^{s}\frac{n_i}{d}\;\tr_{n_i}(x_i),
  \quad
  d =\sum_{i=1}^s n_i^{2}.
\]
Jones–Sunder call $\tau_A$ the \emph{Markov trace} because
$\tau_A\circ E=\tau_A$ for the unique trace-preserving expectation $E$ onto a
unital subalgebra $B$.  Its \emph{trace vector} is precisely $\dimvec n$.

\begin{theorem}[Jones–Watatani conditional expectation]\label{thm:JW-CE}
Let
\[
   A=\bigoplus_{i=1}^{s}M_{n_i}(\mathbb{C}), \qquad 
   B=\bigoplus_{j=1}^{r}M_{m_j}(\mathbb{C})\subseteq A,
\]
and let
\[
  \tau_A\!\bigl((x_1,\dots,x_s)\bigr)=\frac1d\sum_{i=1}^{s} n_i\,\tr_{n_i}(x_i),
  \qquad
  d=\sum_{i=1}^{s}n_i^{2},
\]
be the \textbf{Markov trace}.  
There exists a unique trace-preserving conditional expectation 
\(E_\tau:A\to B\).  
On the $i$-th block \(x_i\in M_{n_i}(\mathbb{C})\) it is given by
\begin{equation}\label{eq:JW-formula}
{\color{MyDarkBlue}
  E_\tau(x_i)=
  \bigoplus_{j=1}^{r}\! 
     \frac{m_j}{n_i}\; 
     \sum_{\ell=1}^{h_{ij}}
        {\bigl(\theta_{ij}^{\,\ell}\bigr)^{-1}\! 
        \bigl(e_{ij}^{\,\ell}x_i\,e_{ij}^{\,\ell}\bigr)}
  \quad(1\le i\le s),
}
\end{equation}
where \(e_{ij}^{\,\ell}:=\theta_{ij}^{\,\ell}(I_{m_j})\) is the supporting projection of the
$\ell$-th copy of \(M_{m_j}\) inside \(M_{n_i}\).
\end{theorem}

\paragraph{Notation at a glance.}
\begin{itemize}
  \item $H=[h_{ij}]$ is the inclusion matrix; $h_{ij}$ counts how many times
        \(M_{m_j}\) appears inside \(M_{n_i}\).
  \item \(\theta_{ij}^{\,\ell}:M_{m_j}(\mathbb{C})\hookrightarrow M_{n_i}(\mathbb{C})\)
        (\(\ell=1,\dots,h_{ij}\)) is the $\ell$-th block-diagonal embedding.
  \item \(e_{ij}^{\,\ell}= \theta_{ij}^{\,\ell}(I_{m_j})\) is the corresponding
        \textit{corner projection}. 
        The compression \(e_{ij}^{\,\ell}x_i e_{ij}^{\,\ell}\) lies in the corner
        algebra \(e_{ij}^{\,\ell}M_{n_i}e_{ij}^{\,\ell}\),
        which the inverse map \((\theta_{ij}^{\,\ell})^{-1}\) identifies with
        \(M_{m_j}\).
  \item The scalar factor \(m_j/n_i\) is chosen so that
        \(\tau_A(x)=\tau_A\!\bigl(E_\tau(x)\bigr)\) for all \(x\in A\).
\end{itemize}

\begin{proof}[Proof of Theorem \ref{thm:JW-CE}]
We split the argument into six transparent steps.

\medskip
\noindent\textbf{Step 1: Block notation and embeddings.} Fix $i$.
Inside the simple block $M_{n_i}(\mathbb{C})$ the inclusion
$B\subseteq A$ picks out
$h_{ij}$ mutually orthogonal corners
\[
   e_{ij}^{\,\ell}= \theta_{ij}^{\,\ell}(I_{m_j}) 
   \quad
   (1\le j\le r,\;1\le \ell\le h_{ij}),
   \qquad
   \sum_{j,\ell}e_{ij}^{\,\ell}=I_{n_i}.
\]
Each $\theta_{ij}^{\,\ell}$ is a unital $\ast$–embedding
\[
\theta_{ij}^{\,\ell}:M_{m_j}\hookrightarrow M_{n_i}
\]
and \(e_{ij}^{\,\ell}M_{n_i}e_{ij}^{\,\ell}\)
inherits the matrix–algebra structure of $M_{m_j}$ through
\((\theta_{ij}^{\,\ell})^{-1}\). 
Explicitly, we may pick matrix units \(\{u^{(\ell)}_{ab}\}_{a,b=1}^{m_j}\subset M_{n_i}\)
so that
$\theta_{ij}^{\,\ell}(E_{ab})=u_{ab}^{(\ell)}$
and
$e_{ij}^{\,\ell}=u^{(\ell)}_{11}+\dots+u^{(\ell)}_{m_jm_j}$.

\medskip
\noindent\textbf{Step 2: Defining $E_\tau$.} Given $x=(x_1,\dots,x_s)\in A$ set
\[
E_\tau(x)=\bigl( E_\tau(x_1),\dots,E_\tau(x_s)\bigr)
\]
with $E_\tau(x_i)$ given by formula~\eqref{eq:JW-formula}.
Because direct sums are orthogonal, it suffices to verify all properties
block–wise, i.e. for a \emph{single} $i$ and $x_i\in M_{n_i}$.
\medskip
\noindent\textbf{Step 3: Well–defined and $B$–bimodular.} For fixed $i$, each term
\[
  (\theta_{ij}^{\,\ell})^{-1}(e_{ij}^{\,\ell}x_ie_{ij}^{\,\ell})
\]
lies in $M_{m_j}(\mathbb{C})$,
hence the orthogonal direct sum
over~$j$ is in $B$.
Further, for $b_1,b_2\in B$ their $i$-th block
commutes with every $e_{ij}^{\,\ell}$ (they act \emph{diagonally} on the same
minimal central projections).  Consequently
\[
   E_\tau(b_1x b_2)=b_1E_\tau(x)b_2,
   \qquad
   x\in A,
\]
so $E_\tau$ is $B$-bimodular.

\medskip
\noindent\textbf{Step 4: Complete positivity.} Each corner map
\[
   C_{ij}^{\,\ell}:x_i\mapsto e_{ij}^{\,\ell}x_ie_{ij}^{\,\ell}
\]
is completely positive (it is a conjugation by a projection).
So is its composition with the $*$–isomorphism
\((\theta_{ij}^{\,\ell})^{-1}\). 
A finite direct sum and a positive scalar multiple of completely positive
maps is completely positive.
Hence $E_\tau$ is \emph{completely} positive.

\medskip
\noindent\textbf{Step 5: Trace preservation.} Compute on a \emph{matrix unit} $u_{ab}^{(\ell)}$ (where $a,b\le m_j$):
\[
   E_\tau(u_{ab}^{(\ell)}) = \frac{m_j}{n_i}\, 
     (\theta_{ij}^{\,\ell})^{-1}
     \!\bigl(u_{aa}^{(\ell)}u_{ab}^{(\ell)}u_{bb}^{(\ell)}\bigr)
   =\frac{m_j}{n_i}\,E_{ab}\in M_{m_j}.
\]
Taking the trace,
\[
\tau_A(u_{ab}^{(\ell)}) = \frac{n_i}{d}\,\delta_{ab},
\]
whereas
\( \tau_A\!\bigl(E_\tau(u_{ab}^{(\ell)})\bigr) = \frac{m_j}{d}\,\delta_{ab}.\)
Because $u^{(\ell)}_{ab}$ sits inside the $j$-th block of $B$ exactly
\(m_j/n_i\) times, the two numbers coincide.
By linearity and faithfulness of $\tau_A$
this equality extends from matrix units to all of $A$:
\[
   \tau_A\circ E_\tau=\tau_A.
\]

\medskip
\noindent\textbf{Step 6: Uniqueness.} Suppose $F:A\to B$ is another conditional expectation with
$\tau_A\circ F=\tau_A$. 
Let $\{e_\mu\}$ be \emph{all} matrix units in every block $M_{n_i}$;
these form a basis of $A$ as a vector space.
For each $\mu$,
\[
\tau_A\!\bigl((F-E_\tau)(e_\mu)^*(F-E_\tau)(e_\mu)\bigr)=0
\]
so $(F-E_\tau)(e_\mu)=0$ by faithfulness.
Hence $F=E_\tau$ on a basis, thus on all of~$A$.
\end{proof}
%--------------------------------------------------------------------


\newpage
%--------------------------------------------------------------------
\subsection{The Jones basic construction for \(B\subseteq A\)}

\paragraph{Standing data.}
Throughout we keep the notation of Theorem~\ref{thm:JW-CE}:
\[
  A=\bigoplus_{i=1}^{s}M_{n_i}(\mathbb{C}), \quad
  B=\bigoplus_{j=1}^{r}M_{m_j}(\mathbb{C})\subseteq A, \quad
  H=[h_{ij}]_{i,j}\in M_{s\times r}(\mathbb{N}_0),
\]
and we work in the GNS Hilbert space
\(L^2(A):=L^2(A,\tau_A)\) where the inner product is
\(\langle x,y\rangle=\tau_A(y^*x)\). 
Left multiplication embeds \(A\) faithfully into \(\mathcal B\!\left(L^2(A)\right)\).

\paragraph{The Jones projection.}
Let \(E_\tau:A\!\to\!B\) be the trace–preserving conditional
expectation from Theorem~\ref{thm:JW-CE}.  
Define
\[
 {\color{MyDarkBlue} e_B : L^2(A)\longrightarrow L^2(A), \quad e_B(\xi)=E_\tau(\xi) \qquad(\xi\in A\subseteq L^2(A)) }
\]
Because \(E_\tau\) is completely positive and unital, $e_B$ extends by
continuity to a bounded operator on all of \(L^2(A)\).

\begin{proposition}[Basic facts about \(e_B\)]\label{prop:Jones-rel}
The operator \(e_B\) satisfies
\( e_B=e_B^{\,*}=e_B^{\,2} \)
and, for every \(x\in A\),
\[
   e_B\,x\,e_B = E_\tau(x)\,e_B
   \quad\text{\emph{(Jones relations)}}.
\]
\end{proposition}

\begin{proof}
\emph{(i) Projection.} For $\xi\in A$ we have
$e_B^{\,2}(\xi)=E_\tau\!\bigl(E_\tau(\xi)\bigr)=E_\tau(\xi)=e_B(\xi)$
because $E_\tau$ is idempotent; adjointness is obvious from
$\langle e_B\xi,\eta\rangle=\langle\xi,e_B\eta\rangle$.

\emph{(ii) Jones relations.} Take $x\in A$ and $\xi\in A$. 
Then
\[
e_B x e_B(\xi)= e_B\!\bigl( xE_\tau(\xi)\bigr)
               = E_\tau\!\bigl(xE_\tau(\xi)\bigr)
               = E_\tau(x)\,E_\tau(\xi)
               = E_\tau(x)\,e_B(\xi),
\]
where we used the $B$–bimodule property of $E_\tau$.
\end{proof}

\paragraph{The first basic construction.}
Set
\[
   A_1 := 
   C^*\!\bigl(A,e_B\bigr) = 
   \operatorname{span}_{\mathrm{C}^*}\{a_1e_Ba_2 : a_1,a_2\in A\}.
\]
By Proposition~\ref{prop:Jones-rel} we obtain a new tower of unital
inclusions
\[
 {\color{MyDarkBlue} B \subseteq A \subseteq A_1 } 
\]
Since $A$ is finite–dimensional, the $C^{*}$–algebra $A_1$ is finite-
dimensional as well.  We now describe its structure block–by–block.

\begin{theorem}[Structure of \(A_1\)]\label{thm:A1-structure}
Relative to the pair \((A,B)\) the inclusion matrix of \(A\subseteq A_1\)
is the \emph{Gram matrix} \(H^{\!\top}H\).  Concretely,
\[
   A_1 \cong 
   \bigoplus_{i=1}^{s}
   M_{n_i}(\mathbb{C})\oplus 
   \bigoplus_{j=1}^{r} h_{ij}^{\,2}\,M_{m_j}(\mathbb{C}).
\]
\end{theorem}

\begin{proof}[Idea of proof]
For each pair $(i,j)$ the corner $e_{ij}^{\,\ell}M_{n_i}e_{ij}^{\,\ell}$ is
$\ast$–isomorphic to $M_{m_j}$.  Products
$a_1e_Ba_2$ with $a_k\in M_{n_i}$ are supported on those corners,
and the Jones relations force exactly $h_{ij}^{\,2}$ copies of
$M_{m_j}$ to appear.  A direct matrix computation
shows that the multiplicity of the $i$–th simple block of $A$
inside $A_1$ is $\sum_{j}h_{ij}^{\,2}$,
i.e. the $(i,i)$-entry of $H^{\!\top}H$.
\end{proof}

\paragraph{Index of the inclusion.}
Set
\[
  [A:B]
  :=\tau_A(e_B)^{-\!1}.
\]
Because $\tau_A(e_B)=1/[A:B]$ is simply the norm–square of the inclusion
matrix, we get
\[
 {\color{MyDarkBlue} [A:B]=\|H\|_2^{\,2}=\sum_{i=1}^{s}\sum_{j=1}^{r}h_{ij}^{\,2}. }
\]
This positive integer is called the \textbf{Jones index}
of the finite–dimensional inclusion \(B\subseteq A\).

\begin{remark}[Key properties at a glance]\leavevmode
\begin{itemize}
  \item
    \(E_\tau\) is the unique $\tau_A$–preserving conditional expectation
    from $A$ onto $B$.
  \item
    \(e_B\) encodes both \(E_\tau\) and the inclusion \(B\subseteq A\).
  \item
    The sequence \(B\subseteq A\subseteq A_1\subseteq A_2\subseteq\cdots\),
    obtained by iterating the construction, carries complete information
    about the inclusion; its abstract combinatorial shadow is the
    \emph{Bratteli diagram} whose first incidence matrix is \(H\).
\end{itemize}
\end{remark}
%--------------------------------------------------------------------


\newpage
\section{Need things from Bachi and Bhat paper}

\begin{proposition}[{\color{MyPurple}Necessary conditions for a unitary orthonormal basis}]
\leavevmode

Let $B=\bigoplus_{j=1}^{r}M_{m_j}(\bbC)\subseteq
      A=\bigoplus_{i=1}^{s}M_{n_i}(\mathbb{C})$ be a unital inclusion with
inclusion matrix $H=[h_{ij}]$, dimension vectors
\[
  \widetilde{\bm m}=(m_1,\dots,m_r)^{\!\mathsf T},\;\widetilde{\bm n}=(n_1,\dots,n_s)^{\!\mathsf T}
\]
and conditional expectation $E\colon A\to B$. 
If $(B\subseteq A,E)$ admits a \textbf{unitary orthonormal basis} of size $d$
(in the sense of Pimsner–Popa), then\vspace{0.4em}

\begin{enumerate}[label=\textup{(\arabic*)}]
  \item \textbf{Spectral condition}
        \[
   H^{\!\mathsf T}\widetilde{\bm n}=d\,\widetilde{\bm m},\qquad
           H^{\!\mathsf T}H\,\widetilde{\bm m}=d\,\widetilde{\bm m},\quad
           HH^{\!\mathsf T}\widetilde{\bm n}=d\,\widetilde{\bm n}.
\]

  \item \textbf{Trace (Markov) condition}\; 
        $E$ preserves the Markov trace
        \[
          \tau_A\!\bigl((x_i)\bigr)=\frac{1}{\sum_{i}n_i^{2}}
          \sum_{i} n_i\,\Tr(x_i),
\]
        and $E$ is the \emph{unique} expectation with this property.

  \item \textbf{Integrality of the norm}
        \[
          d=\lVert H\rVert_2^{\,2}
          =\sum_{i=1}^{s}\sum_{j=1}^{r}h_{ij}^{\,2}\in\mathbb{N}.
\]

  \item \textbf{Quadratic dimension identity}
        \[
          \sum_{i=1}^{s}n_i^{2}
          = 
          d\sum_{j=1}^{r}m_j^{2}.
\]
\end{enumerate}
\end{proposition}


\section{Finite–Dimensional setup \& main results (Bakshi–Bhat)}

\subsection*{Standing data}
\begin{itemize}[leftmargin=2.2em]
  \item $B=\displaystyle\bigoplus_{j=1}^{r}M_{m_j}(\mathbb{C})$,
        $A=\displaystyle\bigoplus_{i=1}^{s}M_{n_i}(\mathbb{C})$
        \quad(finite–dimensional \Cst- or von Neumann algebras).
  \item $H=[h_{ij}]\in M_{s\times r}(\bbN_0)$
        \;(\emph{inclusion matrix}),\quad $H\,\widetilde{\bm m}=\widetilde{\bm n}$ where
        \(\widetilde{\bm m}=(m_j)_j^{\!\mathsf T},\ \widetilde{\bm n}=(n_i)_i^{\!\mathsf T}\).
  \item $E:A\!\to\!B$ a faithful conditional expectation.
  \item \emph{Markov trace} on $A$
        \[
          \tau_A((x_i))=\frac{1}{\sum_i n_i^{2}}\, 
          \sum_{i=1}^{s} n_i\,\Tr(x_i),\qquad
          \text{trace vector }=\widetilde{\bm n}.
\]
  \item A \textbf{unitary orthonormal basis} is a family
        $\{W_\alpha\}_{\alpha=1}^{d}\subset A$
        of unitaries with
        $\;E(W_\alpha^{*}W_\beta)=\delta_{\alpha\beta}\,1_B$
        and
        $\;\sum_{\alpha}W_\alpha E(W_\alpha^{*}X)=X\;(X\in A)$.
\end{itemize}

\subsection*{Core theorem (necessary conditions)}
\begin{theorem}[Bakshi–Bhat; Thm.\,A]\label{thm:BB-necessary}
If such a basis of size $d$ exists, then
\begin{align*}
  &{\color{MyDarkBlue} H^{\!\mathsf T}\widetilde{\bm n}=d\,\widetilde{\bm m}
  \quad\Longleftrightarrow\quad
  H^{\!\mathsf T}H\,\widetilde{\bm m}=d\,\widetilde{\bm m},\; 
  HH^{\!\mathsf T}\widetilde{\bm n}=d\,\widetilde{\bm n};
  \\\[-1pt]
  &\tau_A\circ E=\tau_A\quad\text{(\emph{trace preservation; $E$ is unique})};
  \\\[4pt]
  &d=\lVert H\rVert_2^{\,2}\in\bbN, \quad
  \sum_{i}n_i^{2}=d\sum_{j}m_j^{2}.
\end{align*}
\end{theorem}

\subsection*{Constructions proving sufficiency in special cases}
\begin{theorem}[Explicit bases]\label{thm:BB-sufficient}
\leavevmode
\begin{enumerate}[label=\textup{(\alph*)}]
  \item \textbf{Abelian corner (Theorem B).}\; 
        If $B$ is abelian and conditions in
        Thm.\,\ref{thm:BB-necessary} hold,
        a unitary basis exists
        (quasi–circulant construction $\bigl\{V^{t}U^{t}\bigr\}$). 
  \item \textbf{Simple corner (Theorem C).}\; 
        \begin{itemize}[leftmargin=1.55em]
          \item If $B=M_m(\bbC)$ and $E$ preserves the Markov trace,
                then $(B\subseteq A,E)$ has a unitary basis.
          \item If $A=M_n(\bbC)$ and $E$ preserves the normalized trace,
                then $(B\subseteq A,E)$ has a unitary basis.
        \end{itemize}
  \item \textbf{Weyl-type families.}\; 
        When each non-abelian block of $A$ is $M_n(\bbC)$
        and $B$ is abelian,
        the authors build ``multi–Weyl’’ bases
        extending the classical $(V^{j}U^{k})_{j,k}$ scheme.
\end{enumerate}
\end{theorem}

\subsection{Abelian corner: existence of a unitary orthonormal basis}
\label{sec:abelian-corner}

We prove part~\textup{(a)} of Theorem~\ref{thm:BB-sufficient}.
Throughout this subsection

\[
  B = \bbC^{m_1}\oplus\cdots\oplus\bbC^{m_r}
  \qquad\text{(finite dimensional \textbf{abelian} \Cst–algebra)},
\]

\[
  A =M_{n_1}(\mathbb{C})\oplus\cdots\oplus M_{n_s}(\mathbb{C})
  \qquad\text{(finite dimensional multi–matrix algebra)},
\]
and
\( E : A \twoheadrightarrow B \)
is the given trace–preserving conditional expectation.
Let  

\[
  H =[h_{ij}]_{\,1\le i\le s,\;1\le j\le r}\in M_{s\times r}(\bbN)
\]
be the inclusion matrix (notation from \cite{BakshiBhat2025}).  
The hypotheses of Theorem~\ref{thm:BB-necessary} say precisely that  

\[
  H^\mathsf T H\,\widetilde{\bm m}
  = 
  d\,\widetilde{\bm m},
  \hspace{2em}
  \widetilde{\bm m}=\begin{bmatrix}
    m_1\\[-1pt]\vdots\\[-1pt]m_r
  \end{bmatrix},
  \quad
  d:=\|H\|_2^2\in\bbN,
\]
and that \(E\) preserves the Markov/normalised trace.

\paragraph{Step 1: reduce to a single matrix block.}
Because \(B\) is abelian, each summand \(M_{n_i}(\mathbb{C})\) of \(A\) admits the
\emph{same} inclusion pattern.  It is therefore enough to build a unitary
orthonormal basis (UOB) in the simple corner

\[
  M_{n}(\mathbb{C})\supseteq B_0:=\diag(\bbC,\dots,\bbC)\subset M_n(\bbC),
  \qquad n:=n_1=\cdots=n_s ,
\]
and then replicate it block-wise.  From now on we write
\(A=M_n(\bbC)\) and \(B=B_0\).

\paragraph{Step 2: pick two special unitaries \(U,V\).}
Set \(d=r\) (number of minimal projections in \(B\)).
Let \(\omega:=e^{2\pi i/d}\).
Define  

\[
  U:=\diag(1,\omega,\omega^2,\dots,\omega^{d-1})
      \in M_d(\bbC),
  \quad
  V:=(v_{kl})_{k,l=0}^{d-1},\; v_{kl}:=\frac1{\sqrt d}\,\omega^{kl}
      \in M_d(\bbC).
\]
\begin{itemize}
  \item \(U\) is diagonal and unitary;  
  \item \(V\) is the (\(d\times d\)) discrete Fourier matrix, hence unitary;  
  \item \(VU=\omega\,UV\) (a one-dimensional Weyl relation).
\end{itemize}

\paragraph{Step 3: define the “quasi–circulant” family.}
For each \(t=0,\dots,d-1\) put
\[
  W_t:=V^{\,t}\,U^{\,t}\in M_d(\bbC).
\]
Because \(UV=\omega^{-1}VU\) and \(\omega^d=1\),
 the set \(\{W_t : 0\le t<d\}\) is \emph{closed under multiplication}: \(W_sW_t=\omega^{st}W_{s+t}\).

\paragraph{Step 4: orthonormality with respect to \(E\).}
Identify \(B\) with the diagonal subalgebra of \(A=M_d(\bbC)\). 
The trace–preserving expectation is
\(E(X)=\frac1d\Tr(X)\,I_d\) (\(\Tr=\) ordinary matrix-trace).
For \(s,t\in\{0,\dots,d-1\}\),
\[
  E(W_s^*W_t)
  = 
  \frac1d\Tr\bigl(W_{t-s}\bigr)\,I_d
  = 
  \begin{cases}
    I_d &\text{if }s=t,\\[4pt]
    0   &\text{if }s\ne t,
  \end{cases}
\]
because \(\Tr(W_k)=0\) whenever \(k\not\equiv0\pmod d\).  
Hence the \(W_t\)’s are orthonormal over \(B\).

\paragraph{Step 5: unitarity and spanning.}
Each \(W_t\) is a product of two unitaries, hence unitary.
Since \(M_d(\bbC)\) has dimension \(d^2\) over \(\bbC\) and
\(\{W_t b : 0\le t<d,\;b\in\mathrm{diag}(\mathbb{C}^d)\}\) consists of
\(d\cdot d=d^2\) linearly independent vectors, 
the family \(\{W_t\}_{t=0}^{d-1}\) spans \(A\) as a right \(B\)-module.

\paragraph{Step 6: conclusion for the abelian corner.}
Thus \(\{W_t\}_{t=0}^{d-1}\) is a unitary orthonormal basis of \(A=M_n(\bbC)\)
over its abelian subalgebra \(B\).
Re-instating the direct-sum decomposition of \(A\),
we obtain a UOB for the original inclusion
\(B\subseteq A\). 
\hfill\qedsymbol


\subsection{Simple corner: the case $B=M_m(\bbC)$ \& $E$ Markov}
\label{sec:simple-corner-B}

We turn to part~\textup{(b)} of Theorem~\ref{thm:BB-sufficient}.
First assume

\[
  B =M_m(\bbC),
  \quad
  A =\bigoplus_{i=1}^s M_{n_i}(\mathbb{C}),
  \quad
  E\colon A\twoheadrightarrow B
\]
is the unique $B$–bimodule projection preserving the \emph{Markov trace}
on each block of \(A\).  We will produce an explicit UOB.

\medskip
\noindent\textbf{Step 1: Tensor–product decomposition.} Since \(B\) is simple, as a right–\(B\)–module one checks

\[
  A \cong M_m(\bbC)\otimes\bbC^N, \quad
  N :=\sum_{i=1}^s\!n_i^2.
\]
Under this identification:

\[
  B = M_m(\bbC)\otimes 1_N, \quad
  A = M_m(\bbC)\otimes M_N(\bbC),
\]
and the Markov–trace–preserving \(E\) becomes

\[
  E\bigl(x\otimes y\bigr)
  = 
  \frac1N\,\Tr_N(y)\;\bigl(x\otimes I_N\bigr).
\]

\medskip
\noindent\textbf{Step 2: Unitary basis in the “matrix” factor.} By the abelian–corner case (see \S\ref{sec:abelian-corner}), there exists
a unitary orthonormal basis
\(\{W_t\}_{t=0}^{N-1}\subset M_N(\bbC)\)
for the inclusion \(\bbC^N\subset M_N(\bbC)\)
with respect to the normalized trace on \(M_N\).  Equivalently

\[
  E_N\bigl(W_s^*W_t\bigr)
  =\delta_{st}\,I_N, \quad
  W_t\in M_N(\bbC),
\]
where \(E_N(y)=\tfrac1N\Tr_N(y)\,I_N\).

\medskip
\noindent\textbf{Step 3: Lift to a basis of \(A\) over \(B\).} Define

\[
  U_t := I_m \otimes W_t
  \in M_m(\bbC)\otimes M_N(\bbC)=A, \quad
  t=0,\dots,N-1.
\]
Then each \(U_t\) is unitary in \(A\), and

\[
  E\bigl(U_s^*\,U_t\bigr)
  =E\bigl(I_m\otimes W_s^*W_t\bigr)
  =I_m\otimes E_N(W_s^*W_t)
  =\delta_{st}\,(I_m\otimes I_N),
\]
so \(\{U_t\}\) is orthonormal over \(B\).  Finally, since \(\{W_t\}\)
spans \(M_N\) as a right \(\bbC^N\)–module, \(\{U_t\}\)
spans \(A\) as a right \(B\)–module, giving the desired UOB.

\bigskip
\subsection{Dual case: $A=M_n(\bbC)$ \& $E$ normalized trace}
\label{sec:simple-corner-A}

Now assume

\[
  A=M_n(\bbC),
  \quad
  B\subseteq A
  \quad\text{arbitrary multi–matrix subalgebra}, \quad
  E\colon A\twoheadrightarrow B
\]
preserves the \emph{normalized} matrix–trace \(\tau_n\).  Exactly the same
tensor–product argument applies once one observes

\[
  A\cong\bbC^M\otimes M_n(\bbC), \quad
  B=\bbC^M\otimes I_n,
\]
and one invokes the $B=M_m(\bbC)$ case on the first factor.  Explicitly:

\[
A =\bigl(\bbC^M\bigr)\otimes M_n(\bbC),
 \quad
 B =\bigl(\bbC^M\bigr)\otimes I_n, \quad
 E(f\otimes x)
 =\bigl(E_M(f)\bigr)\otimes x,
\]
where \(E_M\colon \bbC^M\to\bbC^M\) is the normalized–trace expectation.
Pulling back a UOB from \(\bbC^M\subset M_M(\bbC)\) completes the proof.
\hfill\qedsymbol


\subsection{Weyl–type families when $B$ is abelian}
\label{sec:weyl}

Assume

\[
  A = \bigoplus_{i=1}^{s} M_{n_i}(\mathbb{C}), \qquad
  B \cong\bbC^s
  \quad(\text{so $B$ is generated by the minimal central projections }p_i),
\]
and that $E\colon A\to B$ is the trace‐preserving expectation coming from the
standard trace on each matrix summand.  We will exhibit a unitary basis

$$ 
  \bigl\{\,W_{i;\alpha,\beta}:1\le i\le s,\;0\le\alpha,\beta<n_i\bigr\} 
$$ 
of $A$ over~$B$.

\medskip
\noindent\textbf{Step 1: Choose Weyl unitaries on each summand.} For each $i=1,\dots,s$, pick a pair of generating “Weyl” unitaries

\[
  U_i,V_i\in M_{n_i}(\mathbb{C}),
\]
satisfying

\[
  U_i V_i =\omega_i\,V_i U_i,
  \quad
  U_i^{n_i}=V_i^{n_i}=I_{n_i},
  \quad
  \omega_i=e^{2\pi i/n_i}.
\]
Then

\[
  \bigl\{\,U_i^\alpha V_i^\beta
  : 0\le\alpha,\beta<n_i\bigr\}
\]
is a unitary orthonormal basis of $M_{n_i}$ over its centre
(with respect to the normalized trace on $M_{n_i}$).

\medskip
\noindent\textbf{Step 2: Embed these into~$A$.} Under the identification

\[
  A \cong (p_1Ap_1)\oplus\cdots\oplus(p_sAp_s), \quad
  p_iAp_i\cong M_{n_i}(\mathbb{C}),
\]
define for each triple $(i,\alpha,\beta)$ the unitary

\[
  W_{i;\alpha,\beta}
  = 
  \underbrace{0\oplus\cdots\oplus 0}_{i-1}
  \oplus
  \bigl(U_i^\alpha V_i^\beta\bigr)
  \oplus
  \underbrace{0\oplus\cdots\oplus 0}_{s-i}
  \in A.
\]
By construction $W_{i;\alpha,\beta}$ acts nontrivially only on the
$i$–th summand and is hence unitary in~$A$.

\medskip
\noindent\textbf{Step 3: Check orthonormality over~$B$.} Since 

\[
  E\bigl(x_1\oplus\cdots\oplus x_s\bigr)
  =\sum_{i=1}^s\tfrac1{n_i}\,\Tr(x_i)\;p_i,
\]
one computes for any indices $(i,\alpha,\beta),\,(j,\gamma,\delta)$

\[
  E\bigl(W_{i;\alpha,\beta}^* \,W_{j;\gamma,\delta}\bigr)
  =\delta_{ij}\; 
  E_i\!\bigl((U_i^\alpha V_i^\beta)^* (U_i^\gamma V_i^\delta)\bigr)\,p_i
  =\delta_{ij}\,\delta_{\alpha\gamma}\,\delta_{\beta\delta}\;p_i,
\]
where
\( E_i(y)=\tfrac1{n_i}\Tr(y)\,I_{n_i} \)
is the normalized‐trace expectation on the $i$–th block.  In particular

\[
  E\bigl(W_{i;\alpha,\beta}^* \,W_{j;\gamma,\delta}\bigr)
  = \delta_{\,i,j}\,\delta_{\alpha,\gamma}\,\delta_{\beta,\delta}
    \;(p_1+\cdots+p_s)
  =\delta_{(i,\alpha,\beta),(j,\gamma,\delta)}\,1_B.
\]
Hence the collection $\{W_{i;\alpha,\beta}\}$ is orthonormal over~$B$.

\medskip
\noindent\textbf{Step 4: Spanning property.} By counting dimensions one sees there are
\( \sum_i n_i^2=\dim A/\dim B \) 
unitaries above, so orthonormality forces them to span $A$ as a right
$B$–module.  Concretely, any
$x=x_1\oplus\cdots\oplus x_s\in A$
can be written block‐wise by the Weyl‐basis expansion in each
$p_iAp_i$, and hence by the same formulas in $A$ itself.

\medskip
\noindent\textbf{Conclusion.}
The unitaries

\[
  \{\,W_{i;\alpha,\beta}\} = \{\,0\oplus\cdots\oplus U_i^\alpha V_i^\beta\oplus\cdots\oplus0\}
\]
are unitary, orthonormal over $B$, and span $A$.  This completes the proof
of part~\textup{(c)} of Theorem~\ref{thm:BB-sufficient}.
\qed


\newpage
\subsection*{Applications inside the finite–dimensional tower}
\begin{itemize}[leftmargin=1.8em]
  \item \emph{Relative entropy}. 
        If a unitary basis exists then
        \( H(A_1\,|\,A)=\ln\!\bigl(\|H\|_2^{2}\bigr) \)
        where $A_1$ is the first Jones basic construction.
  \item \emph{Depth-2 subfactors}. 
        Depth 2 inclusions with abelian relative commutant
        satisfy the spectral condition automatically,
        hence admit a unitary basis.
\end{itemize}

\subsection*{Author’s conjecture (open problem)}
\begin{center}
\emph{Spectral + trace conditions are \textbf{sufficient} in full generality.}
\end{center}
Current status:
\begin{itemize}[leftmargin=1.8em]
  \item \textbf{Proved} in the special settings of
        Theorem~\ref{thm:BB-sufficient}(a)\,–\,(c).
  \item \textbf{Open} when both $A$ and $B$ are genuinely multi-matrix
        (no simple or abelian side) and $H$ satisfies the PF-eigenvector
        equations but lacks additional regularity.
\end{itemize}

\begin{remark}[Perspective]
The paper bridges classical Hadamard/Weyl unitary bases and
Pimsner–Popa bases in subfactor theory, isolating
spectral–theoretic obstructions and giving positive constructions
in broad finite-dimensional landscapes.
\end{remark}


\newpage
\section{\texorpdfstring{$\bbC\oplus M_{2}(\bbC)\subset M_{3}(\bbC)\oplus M_{9}(\bbC)$}{C ⊕ M₂ ⊂ M₃ ⊕ M₉}: a $2\times2$ inclusion of index $18$}
\label{subsec:2x2-index18}

\begin{example}[Inclusion matrix and basic data]\label{ex:2x2-H}
\leavevmode

\noindent
\textbf{Inclusion matrix.}
\[
  H=\begin{bmatrix}3&0\\
  1&4\end{bmatrix},
  \qquad
  H^{\!\mathsf T}\!H=\begin{bmatrix}10&4\\
  4&16\end{bmatrix}.
\]

\medskip
\noindent
\textbf{Dimension vectors.}
\[
  \dimvec m=\begin{bmatrix}1\\[2pt]2\end{bmatrix},
  \qquad
  \dimvec n=H\dimvec m=\begin{bmatrix}3\\[2pt]9\end{bmatrix}.
\]
Because \(H^{\!\mathsf T}H\,\dimvec m=18\,\dimvec m\), the Perron–Frobenius
eigenvalue is \(d=18\); hence the Jones index satisfies \([A:B]=18\).
\end{example}

\paragraph{Explicit embedding.}
Let \(B=\bbC\oplus M_{2}(\bbC)\) and \(A=M_{3}(\bbC)\oplus M_{9}(\bbC)\). 
For \(\beta\in\bbC\) and \(B_{2}\in M_{2}(\bbC)\)

\[
  \iota\bigl(\beta\oplus B_{2}\bigr)=
  \underbrace{\begin{bmatrix}\beta&0&0\\
  0&\beta&0\\
  0&0&\beta\end{bmatrix}}_{M_{3}(\mathbb{C})}
  \;\oplus\; 
  \underbrace{\begin{bmatrix}
     \beta &         &         &         &         \\
           & B_{2}   &         &         &         \\
           &         & B_{2}   &         &         \\
           &         &         & B_{2}   &         \\
           &         &         &         & B_{2}
  \end{bmatrix}}_{M_{9}(\mathbb{C})\text{ — one 1x1 and four }2\times2\text{ blocks}}.
\]

\begin{proposition}[{\color{MyPurple}Markov traces with the \emph{ordinary} matrix trace}]
\label{prop:Markov-ordinary}
Let

\[
  B=\bbC\oplus M_{2}(\bbC),
  \qquad
  A=M_{3}(\bbC)\oplus M_{9}(\bbC),
\]
and keep the embedding \(\iota:B\hookrightarrow A\) from
\textsection\ref{subsec:2x2-index18}.  With

\[
  m_1=1,\;m_2=2, \qquad
  n_1=3,\;n_2=9, \qquad
  d_B=\sum_j m_j^{2}=5, \quad
  d_A=\sum_i n_i^{2}=90,
\]
define the \textbf{Markov traces}

\[
 {\color{MyDarkBlue} \tau_B(\beta\oplus B_2)
      :=\frac{1}{d_B}
        \bigl(m_1\,\beta+m_2\,\Tr B_2\bigr)
      =\frac{1}{5}\bigl(\beta+2\,\Tr B_2\bigr) }
\]


\[
 {\color{MyDarkBlue} \tau_A(\alpha\oplus A_9)
      :=\frac{1}{d_A}
        \bigl(n_1\,\Tr\alpha+n_2\,\Tr A_9\bigr)
      =\frac{1}{90}\bigl(3\,\Tr\alpha+9\,\Tr A_9\bigr) }
\]
(where \(\Tr\) is the \emph{usual} matrix trace, so \(\Tr(I_n)=n\)).
Then
\begin{enumerate}[label=\textup{(\roman*)}, leftmargin=2.5em]
  \item\label{it:tracial-ord}
        \(\tau_B\) and \(\tau_A\) are faithful traces on \(B\) and \(A\);
  \item\label{it:compat-ord}
        \(\tau_A\circ\iota=\tau_B\) (no extra scale factor);
  \item\label{it:state-ord}
        \(\tau_A(1)=\tau_B(1)=1\); hence both are already \emph{states}.
\end{enumerate}
\end{proposition}

\begin{proof}
\textbf{\ref{it:tracial-ord}.}
Each \(\tau\) is a positive linear combination of ordinary traces, so
traciality and faithfulness are immediate.

\smallskip
\textbf{\ref{it:compat-ord}.}
Take \(x=\beta\oplus B_2\in B\).
Under \(\iota\) it becomes

\[
  \iota(x)=
  \underbrace{\begin{bmatrix}\beta&0&0\\
  0&\beta&0\\
  0&0&\beta\end{bmatrix}}_{ \alpha\in M_3} 
  \;\oplus\; 
  \underbrace{\begin{bmatrix}
     \beta &         &         &         &         \\
           & B_2   &         &         &         \\
           &         & B_2   &         &         \\
           &         &         & B_{2}   &         \\
           &         &         &         & B_{2}
  \end{bmatrix}}_{A_9\in M_9}.
\]
Hence
\( \Tr\alpha=3\beta \)
and
\( \Tr A_{9}=\beta+4\,\Tr B_2. \)
Substituting in \(\tau_A\) gives

\[
  \tau_A\!\bigl(\iota(x)\bigr)
  =\frac{1}{90}\bigl(3\cdot3\beta+9(\beta+4\,\Tr B_2)\bigr)
  =\frac{1}{5}\bigl(\beta+2\,\Tr B_2\bigr)
  =\tau_B(x).
\]

\smallskip
\textbf{\ref{it:state-ord}.}
Apply each trace to the unit \(1_B=(1,I_2)\):

\[
  \tau_B(1)=\frac{1}{5}\bigl(1+2\cdot2\bigr)=1, \qquad
  \tau_A(1)=\frac{1}{90}\bigl(3\cdot3+9\cdot9\bigr)=1.
\]
Thus both traces are normalised states, and no rescaling is necessary.
\end{proof}

\begin{remark}[{\color{MyDarkGreen}Dimension-weighted formula}]\label{rem:dimension-weight}
The coefficients
\( 
  \frac{m_j}{d_B},\frac{n_i}{d_A}
\)
come from the general rule
\( 
  \tau\!\bigl((x_i)\bigr)=\sum_i \dfrac{\dim(V_i)}{\sum_k \dim(V_k)^2}\,\Tr(x_i),
\)
so each simple block contributes in proportion to its
dimension.  With these weights \(\tau_A\) \emph{extends} \(\tau_B\) through
the inclusion map and automatically yields the Jones index
\( [A:B]=\Vert H\Vert_2^2=18. \)
\end{remark}

\paragraph{At-a-glance summary (updated).}
\[
 {\color{MyDarkBlue} B=\bbC\oplus M_{2}(\bbC)
    \overset{\iota}{\subset} 
    A=M_{3}(\bbC)\oplus M_{9}(\bbC) }
\]
\[
  H=\begin{bmatrix}3&0\\
  1&4\end{bmatrix},
  \qquad
  [A:B]=18.
\]

\textbf{Markov traces (states).}

\[
  \tau_B(\beta\oplus B_2)=\frac15\!\bigl(\beta+2\,\Tr B_2\bigr), \qquad
  \tau_A(\alpha\oplus A_9)=\frac1{90}\bigl(3\,\Tr\alpha+9\,\Tr A_9\bigr).
\]

\textbf{Trace-preserving conditional expectation \(E:A\to B\):}
Let $A_9 = (a_{ij})$ and let $B_2^{(\ell)}$ for $\ell=1,..,4$ denote the $\ell$-th $2 \times 2$ block embedded in $M_9$.
\[
  E(\alpha\oplus A_{9})=
  \left[\frac{1}{3}\Tr(\alpha) + \frac{1}{9} a_{11}\right]
  \oplus 
  \frac{2}{9} \sum_{\ell=1}^4 B_2^{(\ell)}(A_9)
\]
where $B_2^{(\ell)}(A_9)$ is the matrix obtained by restricting $A_9$ to the $\ell$-th block.

\begin{figure}[ht]
  \centering
  \missingfigure{Bratteli_for_example.eps}
  \caption[Directed Bratteli diagram for $\bbC\oplus M_{2}\subset M_{3}\oplus M_{9}$] { \small\textbf{Directed Bratteli diagram of the inclusion
\(\bbC\oplus M_{2}(\mathbb{C})\subset M_{3}(\mathbb{C})\oplus M_{9}(\mathbb{C})\).}
Black dots on the left mark the two summands of \(B\):
\(\bbC\) (top, dimension 1) and \(M_{2}(\mathbb{C})\) (bottom, dimension 2).
Black dots on the right mark the summands of \(A\):
\(M_{3}(\mathbb{C})\) (top, dimension 3) and \(M_{9}(\mathbb{C})\) (bottom, dimension 9).
Arrows point from each \(B\)-vertex to each \(A\)-vertex, with multiplicities
given by the inclusion matrix
\(\,H=\bigl[\!\begin{smallmatrix}3&0\\
1&4\end{smallmatrix}\bigr]\): 
three arrows from \(\bbC\) to \(M_{3}\), one arrow from \(\bbC\) to \(M_{9}\),
and four arrows from \(M_{2}\) to \(M_{9}\).
}
  \label{fig:Bratteli-2x2}
\end{figure}


\newpage
\subsection{Path–algebra realisation of the tower}

\begin{definition}[Minimal \emph{vs.} minimal–\emph{central}]\label{def:min-vs-central}
Let \(A\) be a finite–dimensional \Cst-algebra.
\begin{itemize}[leftmargin=2em]
  \item A (nonzero) projection \(p\in A\) is \emph{minimal} if there is no nonzero projection \(q<p\).
  \item A central projection \(z\in Z(A)\) is \emph{minimal central} if no nonzero central \(z'<z\).
  \item Denote \(\pi(A)=\{\text{minimal central projections}\}\) and
        \(\mathsf{Min}(A)=\{\text{all minimal projections}\}\).
\end{itemize}
In the decomposition

\[
  A\cong\bigoplus_{p\in\pi(A)}M_{d_p}(\bbC),
\]
each \(p\) picks out one simple summand, and the ordinary minimal projections
are the rank–one matrix units under a unique \(p\).
\end{definition}

\begin{definition}[Inclusion matrix]\label{def:inc-matrix}
For the unital inclusion \(\cB\subseteq A\) with
\(\cB\cong\bigoplus_{p\in\pi(\cB)}M_{m_p}(\bbC)\) and
\(A\cong\bigoplus_{q\in\pi(A)}M_{n_q}(\bbC)\),
the \emph{inclusion matrix}

\[
  H=\bigl[h_{p,q}\bigr]_{\,p\in\pi(\cB),\,q\in\pi(A)},\quad
  h_{p,q}\in\mathbb{N}_0,
\]
records how many copies of the \(\cB\)-block \(M_{m_p}\) sit inside the
\(A\)-block \(M_{n_q}\). Equivalently,
\(\;H\,\dimvec{m}=\dimvec{n}\), where
\(\dimvec m=(m_p)_{p\in\pi(\cB)}^{\!\mathsf T},\)
\(\dimvec n=(n_q)_{q\in\pi(A)}^{\!\mathsf T}\).
\end{definition}

\paragraph{Edge sets and path spaces.} For each level \(k\ge0\) let

\[
  \Omega_k=\{\text{edges in the Bratteli diagram of }\cB\subseteq A\},
\]
and for \(-1\le\ell<k\) define

\[
  \Omega[\ell,k]
  =\bigl\{(\alpha_{\ell+1},\dots,\alpha_k):\;
     \alpha_i\in\Omega_i,\;\;
     r(\alpha_i)=s(\alpha_{i+1})\; 
  \bigr\},
\quad
  \Omega[\ell,\infty)=\varprojlim_k\Omega[\ell,k].
\]

\paragraph{Hilbert spaces.} Set

\[
  H[\ell,k]:=\ell^2\bigl(\Omega[\ell,k]\bigr), \quad
  H:=H[-1,\infty)=\ell^2\bigl(\Omega\bigr),
\]
with orthonormal basis \(\{\delta_\alpha\}\).

\begin{definition}[Path–algebra blocks]\label{def:scriptBn}
For \(n\ge-1\), define the block

\[
  \mathscr B_n
  :=\Bigl\{\,x\in\mathcal B(H)\;\Big|\;
      \exists\,a\in\mathcal B\bigl(H[n,\infty)\bigr):\;
      x(\alpha,\beta)
      =\delta_{\alpha_{[-1,n-1]},\,\beta_{[-1,n-1]}}\, 
      a\bigl(\alpha_{[n,\infty)},\beta_{[n,\infty)}\bigr)
  \Bigr\}.
\]
Then
\(\mathscr B_{-1}\subset\mathscr B_0\subset\mathscr B_1\subset\cdots\)
is naturally isomorphic (level-by-level) to
\(\cB\subset A_0\subset A_1\subset\cdots\).
\end{definition}

\begin{remark}
A full system of matrix units in \(\mathscr B_n\) is

\[
  e_{\lambda,\mu}(\alpha,\beta)
  =\delta_{\alpha_{[-1,n]},\,\lambda}\; 
  \delta_{\beta_{[-1,n]},\,\mu}\; 
  \delta_{\alpha_{[n+1,\infty)},\,\beta_{[n+1,\infty)}},
\]
for \(\lambda,\mu\in\Omega[-1,n]\) with \(r(\lambda)=r(\mu)\).
\end{remark}

\paragraph{Markov trace on \(\mathscr B_n\).}
Let \(t^{(n)}_q>0\) be the weight on block \(q\in\pi(A_n)\), normalized so
\(\sum_{q}t^{(n)}_q\,n_q=1\).  Then for \(x\in\mathscr B_n\),

\[
 {\color{MyDarkBlue} \tau_{\mathscr B_n}(x)
    = 
    \sum_{\alpha\in\Omega[-1,n]}
      t^{(n)}_{\,r(\alpha)}\; 
      x(\alpha,\alpha) }
\]

\paragraph{Expectation \(E_{\mathscr B_k}:\mathscr B_n\to\mathscr B_k\) \emph{($k<n$)}.}
For \(x\in\mathscr B_n\) and paths \(\alpha,\beta\),

\[
  \bigl(E_{\mathscr B_k}x\bigr)(\alpha,\beta)
  =\delta_{\alpha_{[-1,k]},\,\beta_{[-1,k]}}
  \sum_{\substack{\theta\in\Omega[k,n]\\s(\theta)=r(\alpha_k)}}
    \frac{\,t^{(n)}_{\,r(\theta)}\,}
         {t^{(k)}_{\,r(\alpha_k)}}
    \;x\bigl(\alpha_k\!\cdot\!\theta,\; 
            \beta_k\!\cdot\!\theta\bigr).
\]

\begin{definition}[Jones projection]\label{def:jones-proj}
For \(n\ge0\), the Jones projection
\(e_{n+1}\in\mathcal B(H)\) is defined by

\[
  e_{n+1}(\alpha,\beta)
  =\delta_{\alpha_{[-1,n-1]},\,\beta_{[-1,n-1]}}\, 
  \delta_{\alpha_{[n+1,\infty)},\,\beta_{[n+1,\infty)}}\, 
  \sqrt{\frac{t^{(n)}_{r(\alpha_n)}}{t^{(n-1)}_{s(\alpha_n)}}}\, 
  \delta_{\alpha_n,\,\beta_n}.
\]
\end{definition}


\begin{remark}
\(e_{n+1}\in A_{n-1}'\cap A_{n+1}\), self–adjoint and idempotent, and satisfies

\[
  e_{n+1}\,a\,e_{n+1}
  = 
  E_{A_{n-1}}(a)\,e_{n+1},
  \quad
  a\in A_n.
\]
\end{remark}

\subsection{First Jones basic construction}

Recall

\[
  A_0 = \mathcal{B} = \mathbb{C}\oplus M_2(\mathbb{C}),
  \quad
  A_1 = A = M_3(\bbC)\oplus M_9(\bbC),
  \quad
  H=\begin{bmatrix}3&0\\
  1&4\end{bmatrix}.
\]
The first basic construction produces

\[
  A_2 = C^*(A_1,e_{B}) \cong M_{18}(\mathbb{C})\oplus M_{36}(\mathbb{C}),
\]
because the new inclusion matrix is the transpose
\( H^{T} = \begin{bmatrix}3&1\\
0&4\end{bmatrix}, \)
and hence the dimension vector of \(A_2\) is

\[
  H^{\!T}\,\begin{bmatrix}3\\9\end{bmatrix}
  = 
  \begin{bmatrix}3&1\\
  0&4\end{bmatrix}
  \begin{bmatrix}3\\9\end{bmatrix}
  =\begin{bmatrix}18\\36\end{bmatrix}.
\]
In particular, \(\dim A_2 = 18^2 + 36^2\) and the Perron–Frobenius eigenvalue
of \(H^{T}H\) remains \(d=18\), matching the Jones index \([A:B]=18\).


\section*{Calculation of \( e_2 \) in the Jones Tower}
% Gemini agent note: The following section calculates e_1, but is titled e_2, which is a common convention.

Consider the unital inclusion \( \mathbb{C} \oplus M_2(\mathbb{C}) \subset M_3(\mathbb{C}) \oplus M_9(\mathbb{C}) \), with the tower:
- \( A_0 = \mathcal{B} = \mathbb{C} \oplus M_2(\mathbb{C}) \),
- \( A_1 = A = M_3(\mathbb{C}) \oplus M_9(\mathbb{C}) \),
- \( A_2 = C^*(A_1, e_B) \cong M_{18}(\mathbb{C}) \oplus M_{36}(\mathbb{C}) \).

We compute the Jones projection \( e_2 \) for \( A_1 \subset A_2 \) using the path algebra formulation.
\medskip

\subsection*{Step 1: Bratteli Diagram and Inclusion Matrices}
The inclusion \( A_0 \subset A_1 \) has matrix:

\[
H = \begin{bmatrix} 3 & 0 \\ 1 & 4 \end{bmatrix},
\]
with \( \pi(A_0) = \{ p_1, p_2 \} \), \( \pi(A_1) = \{ q_1, q_2 \} \). For \( A_1 \subset A_2 \), \( \pi(A_2) = \{ r_1, r_2 \} \), and:

\[
H^T = \begin{bmatrix} 3 & 1 \\ 0 & 4 \end{bmatrix}.
\]
\medskip

\subsection*{Step 2: Path Algebra Formula}
For \( e_2 \) (corresponding to the inclusion $A_0 \subset A_1$, so $n=0$ in the formula for $e_{n+1}$):

\[
e_{1}(\alpha, \beta) = \delta_{\alpha_{[-1,-1]}, \beta_{[-1,-1]}} \, \delta_{\alpha_{[1, \infty)}, \beta_{[1, \infty)}} \, \sqrt{ \frac{t^{(0)}_{r(\alpha_0)}}{t^{(-1)}_{s(\alpha_0)}} } \, \delta_{\alpha_{0}, \beta_{0}}
\]
where for \( \alpha, \beta \in \Omega[-1, \infty) \):
- \( \alpha = (\alpha_0, \alpha_1, \dots) \) where \( \alpha_i \) is an edge from level $i-1$ to $i$.
- \( e_1(\alpha, \beta) \neq 0 \) only if the paths agree after level 1 and have the same first edge $\alpha_0=\beta_0$.
\medskip

\subsection*{Step 3: Weights}
The Markov trace weights for the vertices (simple blocks) are proportional to the dimensions of the corresponding simple modules. Let's normalize $t^{(-1)}=1$ for the unique block at level -1.
\begin{itemize}
    \item Level 0: \( \dimvec{m} = (1, 2)^T \). The trace vector is \( t^{(0)} = (1, 2)^T \).
    \item Level 1: \( \dimvec{n} = (3, 9)^T \). The trace vector is \( t^{(1)} = (3, 9)^T \).
\end{itemize}
The coefficients for the square root term are:
\begin{itemize}
    \item Path from level -1 to $p_1$: \( \sqrt{ \frac{t^{(0)}_{p_1}}{t^{(-1)}} } = \sqrt{ \frac{1}{1} } = 1 \).
    \item Path from level -1 to $p_2$: \( \sqrt{ \frac{t^{(0)}_{p_2}}{t^{(-1)}} } = \sqrt{ \frac{2}{1} } = \sqrt{2} \).
\end{itemize}
\medskip

\subsection*{Step 4: Matrix Form}
The trace of the Jones projection must be \( \tau_{A_2}(e_2) = 1/[A_1:A_0] = 1/18 \). The Markov trace on $A_2$ is \( \tau_{A_2}(X_1 \oplus X_2) = \frac{1}{1620}(18\Tr(X_1) + 36\Tr(X_2)) = \frac{1}{90}\Tr(X_1) + \frac{1}{45}\Tr(X_2) \). Let \( e_2 = P_1 \oplus P_2 \), where $P_1, P_2$ are projections. Their ranks must satisfy \( \frac{1}{90}\mathrm{rank}(P_1) + \frac{1}{45}\mathrm{rank}(P_2) = \frac{1}{18} \), which simplifies to \( \mathrm{rank}(P_1) + 2\,\mathrm{rank}(P_2) = 5 \).
A possible solution for the ranks is \( (\mathrm{rank}(P_1), \mathrm{rank}(P_2)) = (3, 1) \). This leads to the representation:

\[
e_{2} = \begin{bmatrix} I_3 & 0 \\ 0 & 0 \end{bmatrix}_{18 \times 18} \oplus \begin{bmatrix} 1 & 0 \\ 0 & 0 \end{bmatrix}_{36 \times 36}.
\]


%──────────────────────────────────────────────────────────────────────────────%
%                               BIBLIOGRAPHY                                  %
%──────────────────────────────────────────────────────────────────────────────%
\begin{thebibliography}{9}

\bibitem{JonesSunder}
 V.F.R.~Jones and V.S.~Sunder,
 \emph{Introduction to subfactors},
 Cambridge University Press, 1997.

\bibitem{BakshiBhat2025}
 K.C.~Bakshi and B.V.R.~Bhat,
 \emph{Unitary orthonormal bases of finite–dimensional inclusions},
 preprint, 2025.
\end{thebibliography}
\end{document}